Para que as operações de encontro e atracação \textit{(rendezvous and docking/berthing)} entre espaçonaves possam ser realizadas com sucesso, é necessário definir os sistemas de referência envolvidos. Um sistema de referência é composto por uma origem e três eixos ortogonais e, estes descrevem a posição e a orientação de corpos no espaço tridimensional.

Neste capítulo, serão abordados os principais sistemas de referência utilizados nas manobras de rendezvous and docking/berthing (RVD/B), especialmente utilizando os sistemas de referência inercial e o LVLH (\textit{Local Vertical Local Horizontal}). Para descrever e compreender os movimentos das espaçonaves corretamente, se faz necessário obter a correta definição desses quadros, pois permite o planejamento e a execução precisa das manobras orbitais.

Além disso, serão introduzidos conceitos essenciais de cinemática e dinâmica de corpos rígidos aplicados a veículos espaciais. A modelagem matemática desses fenômenos inclui a formulação das Equações de Hill (ou Equações de Clohessy-Wiltshire), que descrevem o movimento relativo entre dois corpos em órbita, além de representações da orientação utilizando ângulos de Euler e quaternions — ferramentas cruciais para o controle de atitude durante a fase de aproximação e acoplamento.

\section{Sistemas de Referência Inerciais e Orbitais}

Em geral, três sistemas de referência são empregados para descrever o movimento de uma espaçonave em operações de encontro:

\begin{enumerate}[label=(\alph*)]
    \item \textbf{Referencial inercial (ECI – Earth-Centered Inertial):} é um sistema de coordenadas cartesiano com origem no centro da Terra e seus eixos estão fixados em relação às estrelas distantes, desconsiderando a rotação da Terra. O \textit{eixo X} aponta na direção do equinócio vernal (\textit{vernal equinox}), o \textit{eixo Z} é perpendicular ao plano equatorial terrestre, apontando para o Polo Norte celeste, e o \textit{eixo Y} é definido de modo a completar um sistema ortonormal de mão direita. Este referencial é utilizado para descrever o movimento orbital de satélites e outros corpos em relação à Terra, especialmente em análises dinâmicas e de propagação orbital.
    
    \item \textbf{Referencial orbital (LVLH – Local Vertical Local Horizontal):} é um sistema de referência local centrado na espaçonave-alvo \textit{(target)}. O \textit{eixo z} aponta na direção radial, do centro de massa da espaçonave para o centro da Terra; o \textit{eixo x} aponta na direção do movimento orbital (tangente à órbita, na mesma direção do vetor de velocidade); e o \textit{eixo y} é perpendicular ao plano orbital, orientado de modo a completar um sistema ortonormal de mão direita. Este referencial é utilizado para descrever o movimento relativo de um veículo de aproximação \textit{(chaser)} em relação à espaçonave-alvo.
    
    \item \textbf{Referencial de atitude da espaçonave:} este frame é centrado no centro de massa da espaçonave e fixo à sua estrutura. Sua função é descrever a orientação e as manobras de controle de atitude da espaçonave. Os eixos são definidos como: o \textit{eixo x} aponta no sentido da frente da espaçonave, o \textit{eixo y} aponta para a lateral direita (em relação a quem observa a espaçonave, olhando para frente), e o \textit{eixo z} aponta para baixo.
\end{enumerate}

A escolha apropriada dos sistemas de referência permite o uso de modelos matemáticos simplificados para descrever o movimento relativo, como é o caso das Equações de Hill, válidas sob a hipótese de órbitas circulares e pequenas separações entre as espaçonaves.